\magnification=1200

\baselineskip=\normalbaselineskip
\voffset=2\normalbaselineskip
\vsize=8.6truein
\topskip=24truept
\noindent
\headline={\hfill \vbox{ \baselineskip=\normalbaselineskip
\halign{#\hfil\cr
David Koski\cr
\cr
\cr}}}

%\input /u/dkoski/hw/verbatim.tex
				% import verbatim mode

%\input epsf			% allow pictures to be imported
%\epsffile{foo.ps} 		% import the file
%\def\epsfsize#1#2{.5#1}	% scaling


\vskip 30truept

\centerline {\bf COMMON LIBRARY}

\beginsection 1. Overview

% what does it do
% high level design
% other goals

The main goal of the Common library is to provide routines for message passing.
As a side effect, it also gives memory management, queue, and debugging
functions.

\beginsection 2. Design

% more detailed design
% description of processes, data structures, modules

The common library normally uses three processes, or five if {\tt send\_debug()}
is used.  One process is the initial process that calls {\tt main()}:  it's
last step is to call {\tt control()}.  
The other two standard processes are Reader and Writer.  These processes
handle all of the inter-processor communication.  Their specifics are discussed
below.

If a call to {\tt send\_debug()} is made, two processes, DebugBuf and DebugSend
are allocated and run.  This is transparent to the user, and it should not
be a problem.

The common library provides two data structures: {\tt command\_form} and
{\tt queue}.  The {\tt command\_form} is used as an array of function
pointers.  The {\tt queue} data structure is a simple integer queue.  This
is used in many places to provide buffered I/O.

The common library consists of 7 modules:
\halign{\qquad\tt#\hfil&\qquad\vtop{\noindent\hsize=4.3in#\vskip3truept}\hfil\cr
control& The control module defines the {\tt control()} function.  This simply
	waits on the {\tt Writer\_to\_Control} channel for a message.  The
	messages are dispatched according to the {\tt command} structure.
	This module has one file, {\tt control.c}.\cr
var& The var module consists of two files {\tt var.c} and {\tt var\_sr.c}.
	{\tt var()} function handles all {\tt READ\_VAR} and {\tt SET\_VAR}
	messages.  The {\tt read\_var()} and {\tt set\_var()} functions are
	the standard methods for setting variables.\cr
debug& The debug module is one file, {\tt debug.c}, and it provides one
	function: {\tt send\_debug(char *fmt, ...)}.  This function allows
	remote nodes to send debugging messages (like {\tt printf}) up to the
	Sun for display.  This function uses a separate process and channel
	to communicate with the Reader and Writer, so it is safe to call from
	any function.  Additionally, it has internal locking, so concurrent
	calls will not pose a problem.\cr
queue& The queue module is made up of one file, {\tt queue.c}, and it
	implements a simple integer queue.  The functions that are provided
	are:  {\tt queue *q\_alloc()}, {\tt enqueue(q, e)}, {\tt int deq(q)},
	{\tt int emptyq(q)}, and {\tt int fullq(q)}.  No checking is done for
	overflow, the user of the routines must check if that is possible.
	Also, no locking is done in the queue module, if concurrent access is
	possible, it would be wise to place semaphores around the calls.\cr
mem& This is a memory allocation module that has one file, {\tt mem.c}.  The
	memory is allocated using a table showing which blocks are allocated.
	This is needed because we have to be able to free pieces of memory if
	they are in the heap (this heap, not the heap provided by libc) and
	be able to control exactly where the memory is allocated from.  This
	may not be the fastest way to do it (in fact, I can guarantee it is
	not, but it is fairly simple and serves the purpose).  This module
	provides three routines: {\tt mem\_init()}, {\tt mem\_alloc()}, and {\tt
	mem\_free()}.  The memory to be allocated from is defined in the {\tt
	commands.c} file using the {\tt \_MEMORY(size);} macro and defining the
	pointer {\tt char *memory = (char *) address}.  Memory is always
	allocated in MEM\_BLOCK\_SIZE chunks.\cr
talk& This module defines the Reader and Writer processes and is one file, 
	{\tt talk.c}.  The Reader process handles all incoming channels (local
	and physical), and the Writer passes the messages along via a local or
	physical channel.  The Reader process watches the list of channels
	defined by {\tt chan\_list[]}.  If a message can be read (ie it fits
	in the queue of read messages and the channel that it comes from has
	not used up its fair share already), Reader reads the message and
	places it in a message queue.  The Writer signals that it is ready to
	for another message to write by writing to the {\tt Reader\_to\_Writer}
	channel.  If a message is available, the Reader passes one along.
	Messages are guaranteed to be sent along in the order that they
	were read.  If the message is for a different node, the {\tt
	router()} function is called to provide a channel to write to,
	otherwise the message is sent on the {\tt Writer\_to\_Control}
	channel to be handled locally.  The {\tt SEND(header, buffer,
	channel)} macro is provided so that processes can send messages to
	the Reader.\cr
proc& This module is one file, {\tt proc.c} and it defines two routines and a
	variable that are needed for downloading code on the fly (the
	sequencer, for example).  The {\tt prog} variable should be placed in
	the {\tt var\_ptr} structure in the {\tt commands.c} file, and the
	{\tt kill\_proc()} and {\tt execute\_proc()} functions should be in
	the {\tt commands} structure.\cr 
}

The following channels are used by the common library:
\halign{\qquad\tt#\hfil&\qquad\vtop{\noindent\hsize=3.55in#\vskip3truept}\hfil\cr
PROC\_to\_Control& This is a channel that is used by downloaded code
	(microcode) to signal that it has completed (ie when it is killed or
	voluntarily dies).  This is necessary because there is no way to
	actually kill another process on the transputer.  The {\tt cntrl\_reg}
	is set to {\tt DIE}, and it is up to the microcode to {\tt suicide()}.
	A single integer is sent over this channel.  Note:  this is only used
	if downloaded code is being used (by including {\tt \$(OPTIONAL\_OBJS)}
	in the {\tt \$(OBJS)} line.\cr
chan\_list[]& This is a {\tt NULL} terminated list of channels for the Reader
	to watch for messages.  The hardware links, {\tt DEBUG\_CHAN}, {\tt
	Control\_to\_Reader}, and {\tt Reader\_to\_Writer} (as the last channel)
	must be in this list.  The standard protocol should be used over these
	channels.\cr
out\_chan[]& This is a list of 4 channels that the transputer can use to
	communicate between transputers.  Normally these are set to {\tt
	LINKxOUT}.\cr
hw\_chan[]& This is a list of 4 channels that are external input channels.
	This lets the Reader know which channels have non local messages.
	These are normally set to {\tt LINKxIN}.\cr
DEBUG\_CHAN& This is the connection between Debug and Talk, the standard
	protocol is used to send debugging messages up to the sun.  The
	DebugSend process has ownership of this channel.\cr
Proc\_to\_Debug& {\tt send\_debug()} uses this channel to send a header and a
	buffer pointer off to the DebugBuf process.\cr
Buf\_to\_Send& DebugSend writes an integer 0 to this channel to signal that it
	is ready for another message, and DebugBuf writes a header and buffer
	back.\cr
Reader\_to\_Writer& The Writer process writes an integer 0 to this link to
	signal that it is ready for another message, and the Reader writes a
	header and buffer back when one is available.\cr
Writer\_to\_Control& This is the channel that the Writer sends all messages
	that have the local transputer as the destination through.  This
	channel obeys the standard protocol.\cr
Control\_to\_Reader& This is a channel from the Control process back to the
	Reader, and it follows the standard protocol.\cr
}

The global variables that the common library uses are listed in the {\bf
Interfaces} section.

\beginsection 3. Files

% file descriptions
\halign{\qquad\tt#\hfil&\qquad\vtop{\noindent#\vskip3truept}\hfil\cr
Makefile& Makefile for the project.\cr
README& A description of the directory and the files.\cr
README.nodes& Notes on how the nodes are numbered.\cr
common.h& The common library header file.\cr
common.tex& The documentation for the routines in the common library.\cr
control.c& Top level control function.\cr
debug.c& Debugging routines.\cr
mem.c& Memory allocation routines.\cr
message.c& Some message handling stubs.\cr
monitor.c& A channel monitor, not used, never tested.\cr
proc.c& Support for downloaded code.\cr
protocol.h& The message protocol header file.\cr
queue.c& Queue manipulation routines.\cr
shared.h& Shared memory (between the B016 and Sun) header file.\cr
talk.c& Message passing processes, plus a router.\cr
var.c& This module handles the set/read var commands.\cr
var\_sr.c& This is the low level set/read command.\cr
}

\beginsection 4. Interfaces

To use the common library, simply place {\tt include ../com.make} in
your makefile and tack {\tt \$(COMMON\_OBJS) \$(OPTIONAL\_OBJS)} onto
the end of your {\tt OBJS} line.
The file {\tt commands.c} should be in each transputer directory.  It defines
how the common library should act on messages.
This file should include {\tt <conc.h>} and {\tt "common.h"}
in that order.  These are the things that must be defined in {\tt
commands.c}:

\halign{\qquad\tt#\hfil&\qquad\vtop{\noindent\hsize=4in#\vskip3truept}\hfil\cr
MEMORY&
		This tells the memory management routines how much memory is
		available.  This is just a macro which defines some variables,
		and it wants the number of bytes as an argument.  An example
		is {\tt MEMORY(3*1024*1024 + 512*1024);} to reserve 3.5
		megabytes of memory.\cr
command&
		Is an array of type {\tt command\_form} which contains pointers
		to functions.  These functions correspond to the protocol
		commands, one function for each command.  If a command
		is to be ignored you can use a {\tt NULL} function pointer
		to have the controller dispose of the message.\cr
num\_commands&
		Tells the controller how many positions there are in the command
		array.  This is normally 18 for all messages.\cr
var\_ptr&
		Is an array of integer pointers which correspond to the
		protocol variables.\cr
setv\_cmnds&
		Is an array of functions pointers (type {\tt command\_form})
		which match up with the variables in the {\tt var\_ptr} array.
		These functions are called when a variable is being set.
		These can all be set to {\tt set\_var}.\cr
readv\_cmnds&
		Is an array of functions pointers (see above) for actions to
		take when reading a variable.  These can be set to 
		{\tt read\_var}.\cr
chan\_list&
		Is an array of channels for the controller to watch.  Because
		you will probably use dynamically allocated channels, you will
		have to put {\tt NULL} placeholders in and fill it in at 
		runtime.\cr
}

\vskip 12truept
\noindent
Refer to the {\tt commands.c} files currently in the transputer directories
for good examples.

\vskip 12truept
\noindent{\tt main()} also has several things to do:

{\obeylines\obeyspaces\tt
        mem\_init();                            /* initialize memory */
\vskip12truept 
        Reader\_to\_Writer = ChanAlloc();        /* allocate the channels */
        Control\_to\_Reader = ChanAlloc();
        Writer\_to\_Control = ChanAlloc();
\vskip12truept 
        Preader = ProcAlloc(reader, 4096, 0);  /* allocate the processes */
        Pwriter = ProcAlloc(writer, 4096, 0);
\vskip12truept 
        ProcRun(Preader);                      /* run the processes */
        ProcRun(Pwriter);
\vskip12truept 
        control();                             /* start the library */
}

The final interface is the standard protocol (defined in {\tt protocol.h}).
This protocol defines a message format for communicating between different
transputers.  Each message consists of a 4 byte header, followed by 0 or more
integers of data.

The header consists of four 8 bit fields that give the destination node, the
length of the message in integers (not including header), the message type,
and an optional from node.  These fields can be set using the {\tt
TO\_NODE(n)}, {\tt OF\_LENGTH(l)}, {\tt MESSAGE(m)}, and {\tt FROM\_NODE(n)}
macros.  The values can be or'd together to produce a header.  The values can
be extracted from the header with the {\tt NODE(h)}, {\tt LENGTH(h)}, {\tt
MESSAGE(h)}, and {\tt FROM\_WHERE(h)} macros.

These are the currently defined messages (illegal messages are ignored):

\halign{\qquad\tt#\hfil&\qquad\vtop{\noindent\hsize=4in#\vskip3truept}\hfil\cr
IMAGE\_DONE& This message is sent from the dspw's to the b016 when an image is
	ready to be uploaded.  The dspw can send as many of these as it wants
	(assuming it has the pictures ready, of course), and the b016 will
	tell it to send a picture when it is ready.  This message has no
	data.\cr
BEGIN\_XMIT& This message has two formats and meanings.  The first is when the
	b016 sends this message to a dspw.  The message means that the b016 is
	ready for the dspw to upload a picture.  This message has one integer
	argument, the block size.  All uploaded pieces of data must be a
	multiple of the block size in length.  The second use of this message
	is when the dspw is sending a piece of the picture to the b016.  The
	message has data that is some multiple of the block size in length.\cr
SET\_VAR& This message is used to set a variable on a remote transputer.  The
	available variables are described with each transputer.  The message
	arguments have two pieces:  the variable number (and optional offset),
	and the data to set.  The variable number and offset are encoded into
	a single integer: the offset (in integers) is in the upper 16 bits,
	and the variable index is in the lower 16 bits.  This message is
	usually handled by {\tt set\_var()} is {\tt var\_sr.c}, but if it
	is not, the message may have a different format.\cr
READ\_VAR& This message is a request to read a variable off of a remote
	transputer.  This message has two required arguments:  the variable
	to read (encoded in just like {\tt SET\_VAR} requires), and the
	node to send the reply to.  This message also takes an optional
	third argument: length.  The default length is one integer, but
	if zero is given, a string is sent or the requested number of
	integers.  The reply message is a {\tt VAR\_READ}.\cr
DESBUG\_MSG& This is used to send a text message up to the Sun.  The data
	with this message is just the string in transputer format (byte
	swapped!).  These messages can be generated using {\tt send\_debug()}.
	\cr
ABORT\_MSG& Not used.  Aborting is done differently.\cr
STOP\_MSG& Not used.  This is the same as abort, except it will save the
	results so far.\cr
START\_MSG& This message should be sent to the dspw's when picture taking
	is about to start.  This message has no extra data.\cr
PAUSE& Not used.  This would pause the data taking.\cr
RESUME& Not used.  This would resume the data taking.\cr
VAR\_READ& This message is the reply from {\tt READ\_VAR}.  There are two
	fixed size fields:  variable number, and the node it is from.  The
	requested data follows.\cr
READ\_HK& The instrument controller responds to this message.  This is a
	request to read the latest housekeeping data.  The instrument
	controller responds by setting the {\tt hk\_data} variable on the b016,
	which causes the {\tt hkserv} program to display the values.\cr
SET\_VAR\_AKK& This message works just like {\tt SET\_VAR}, except that it
	must send a {\tt AKK} or {\tt AKK\_FAIL} to inform the setter
	(the Sun) that it succeeded or not.\cr
AKK& This message works just like the {\tt DEBUG\_MSG}, plus it also
	means that a {\tt SET\_VAR\_AKK} succeeded.\cr
AKK\_FAIL& This message works just like the {\tt DEBUG\_MSG}, plus it also   
	means that a {\tt SET\_VAR\_AKK} failed.\cr
KILL\_PROC& This message tells the transputer to kill the currently running
	downloaded procedure.  When the process has died (for more info, see
	the {PROC\_to\_Control} entry above), the upper level (tcl) wants to
	see ``Stopped.'' This message should be sent even if no process was
	running.\cr
EXECUTE\_PROC& This tells the transputer to start a new process that runs the
	code that was downloaded into the {\tt prog[]} array.  The new process
	is run at high priority with a stack of size {\tt P\_STACK}.\cr
CHECK\_HK& This message is only recognized by the instrument controller, and
	it takes no data.  The instrument controller will respond by sending
	{\tt DEBUG\_MSG}'s up to the Sun saying either ``{*** %d VOLTAGES
	OUT OF SPEC ***}'' or ``{All hardware-determined voltages are
	within tolerance.}''\cr
ERROR& Not used.  This is to send error messages.\cr
}

Note:  the format of these messages is not strictly enforced.  Some
transputers (the instrument controller in particular) require a slightly
different format.

\beginsection 5. Troubleshooting

Since this library is used on all of the transputers, it is believed to be
thoroughly debugged.  Since this library provied the means to debug a remote
transputer (ie message passing, debugging messages, and memory allocation),
it it is not working there is a big problem.  The best way to track things
down is to get the failure to occur on the b016 and then follow the steps
outlined in the b016 troubleshooting section.

\beginsection 6. Notes

None.
\end
